\documentclass[11pt]{article}
\usepackage[utf8]{inputenc}
\usepackage[T1]{fontenc}
\usepackage{amsmath,amssymb}
\usepackage{booktabs}
\usepackage{graphicx}
\usepackage{multirow}
\usepackage{xcolor}
\usepackage{hyperref}
\usepackage[margin=1in]{geometry}
\usepackage{natbib}
\usepackage{tablefootnote}

\title{Structural Features Provide Strong Baselines for Agent Progress Stagnation Detection:\\A Gold-Validated Empirical Study with Weak Supervision}

\author{Anonymous}
\date{}

\begin{document}
\maketitle

% ============================================================
\begin{abstract}
We present a systematic empirical study of progress stagnation detection in LLM coding agents through weak supervision.
Three labeling functions (LFs)---command repetition, file diff frequency, and outcome backpropagation---generate noisy window-level labels that are aggregated via majority vote.
Our central finding, validated on a 650-window human-annotated gold set, is that stagnation signals are fundamentally structural: 30 handcrafted features achieve 0.84 gold macro-F1 in-domain and an estimated 0.82 cross-agent, surpassing neural embedding probes (0.79--0.81) and an LLM-as-judge baseline (0.71).
However, circular evaluation---where models trained on LF-derived labels are evaluated against those same labels---inflates reported performance by up to 15.5 percentage points for structural methods.
We introduce a hybrid ensemble (structural + probe) that achieves the best gold-validated performance (0.86 macro-F1) by leveraging complementary semantic signals in regions where LF coverage is weakest.
We further characterize LF temporal dynamics: command repetition leads stagnation onset by ${\sim}1.87$ steps on SWE-Agent but reverses to a lagging indicator on OpenHands, where its removal improves F1 by 21.6pp.
Despite these well-motivated temporal patterns, three temporal label model variants all fail to improve over majority vote---a negative finding that constrains future research on sequential weak supervision.
Our results support lightweight structural monitoring as the primary detector, with semantic probes providing targeted complementary value, and highlight circular evaluation as a methodological pitfall for the weak supervision community.
\end{abstract}

% ============================================================
\section{Introduction}
\label{sec:intro}

LLM-based coding agents such as SWE-Agent \citep{yang2024swe} and OpenHands \citep{wang2024openhands} can autonomously resolve software engineering tasks, yet they frequently enter periods of \emph{progress stagnation}---sequences of actions that appear locally reasonable but collectively fail to advance toward the goal.
Unlike simple repetition loops detectable by cycle checkers, semantic stagnation involves diverse commands and file modifications that mask the lack of genuine progress.
Recent benchmarking confirms this is pervasive: SlopCodeBench \citep{slopcodebench2026} reports verbosity increases in 90\% of long trajectories and quality erosion in 80\%.

Existing approaches to agent progress assessment fall into several categories, none fully addressing the monitoring use case.
Process Reward Models (PRMs) \citep{agentprm2026,prm2025,toolprmbench2026} provide step-level training signals but require outcome labels and RL pipelines.
Self-reflection methods \citep{shinn2023reflexion,yao2024tree} rely on the stuck agent's own judgment.
Cycle detectors \citep{SpecRA,TrajAD} catch structural repetition but miss semantically diverse stagnation.
Sequence-level monitoring has been advocated for safety \citep{iclr2025monitoring,weaktostrong2025} but not systematically studied for progress detection.

We frame stagnation detection as a window-level binary classification problem and study it through weak supervision.
Three labeling functions (LFs) generate noisy labels from observable signals; majority vote aggregates them.
We then compare four detection paradigms: (1)~structural feature classifiers, (2)~neural embedding probes, (3)~LLM-as-judge, and (4)~a hybrid ensemble.

Our study makes five contributions:
\begin{enumerate}
    \item \textbf{Structural baseline discovery.} Structural features provide a strong, efficient, and transferable baseline for stagnation detection, achieving 0.84 gold-validated macro-F1 in-domain and 0.82 estimated cross-agent---decisively surpassing probes (0.79) and LLM judge (0.71).
    \item \textbf{Hybrid complementarity.} Semantic probes contribute complementary signals in LF-weak regions (conflict windows, cmd\_rep-only stagnation, high command diversity stagnation), yielding a hybrid ensemble at 0.86 gold macro-F1.
    \item \textbf{Circular evaluation analysis.} We quantify a ``circularity gradient'' where methods closer to LF inputs suffer greater evaluation inflation (structural: 15.5pp, probe: 11.5pp, LLM judge: $-$4.8pp), offering methodological guidance for weak supervision evaluation.
    \item \textbf{LF temporal dynamics.} We characterize systematic, agent-specific temporal roles of LFs and demonstrate that command repetition's role reversal across agents has actionable implications for LF design.
    \item \textbf{Temporal label model negative finding.} Three temporal label model implementations (TWV, HMM, lag-aware HMM) all fail to improve over majority vote, with TWV degrading downstream classifiers by up to 21.7pp.
\end{enumerate}

% ============================================================
\section{Problem Formulation}
\label{sec:method}

\subsection{Window-Level Stagnation Detection}

Given an agent trajectory $\tau = (a_1, a_2, \ldots, a_T)$ of actions, we partition $\tau$ into overlapping windows of size $k$ (default $k{=}5$).
Each window $w_t = (a_{t}, \ldots, a_{t+k-1})$ is labeled as \textsc{Progress} or \textsc{Stagnation}.

\subsection{Labeling Functions}

Three LFs vote independently on each window:

\paragraph{Command Repetition (cmd\_rep).} Flags stagnation when the ratio of repeated commands within a window exceeds a threshold. Fires at 97.4\% of windows on SWE-Agent, making it the dominant signal. It leads stagnation onset by ${\sim}1.87$ steps on SWE-Agent but reverses to a lagging indicator (+3.32 steps) on OpenHands.

\paragraph{File Diff Frequency (file\_diff).} Measures the rate of meaningful file modifications. The most objective signal, with 0.88 human agreement. Leads by $-$4.9 steps on OpenHands and $-$1.2 steps on SWE-Agent.

\paragraph{Outcome Backpropagation (OB).} Propagates task outcome labels backward through the trajectory, tagging late windows of failed tasks as stagnation. Lags by +6.7 steps (SWE-Agent) and +18.1 steps (OpenHands). Removing OB drops stagnation F1 by 22pp, confirming its unique signal contribution despite its privileged nature.

\paragraph{Aggregation.} Majority vote (MV) over the three LFs produces final labels, with a confidence score $\text{conf} \in \{1, 2, 3\}$ reflecting the number of agreeing LFs.

\subsection{Detection Methods}

\paragraph{Structural Features (30 features).} Handcrafted features including command repetition ratio, unique command ratio, file diff count, action diversity, edit distance statistics, and temporal trend features. Trained with Random Forest (RF) and Gradient Boosting (GB).

\paragraph{Neural Embedding Probes.} Three encoder backbones---Qwen3-Embedding-0.6B, BGE-m3, and E5-large-v2---produce window representations via mean pooling of per-step embeddings. A linear probe (sigmoid output) is trained on LF-derived labels.

\paragraph{LLM-as-Judge.} Claude Sonnet performs zero-shot binary classification from the raw action sequence, with no exposure to LF labels.

\paragraph{Hybrid Ensemble.} Probability averaging of RF calibrated probabilities and probe sigmoid probabilities.

\subsection{Gold Set Annotation}
\label{sec:gold}

To break circular evaluation, we construct a 650-window gold set via stratified sampling (50 conf=3, 300 conf=2, 300 conf=1) with dual annotation and adjudication.
Inter-annotator agreement is $\kappa{=}0.73$ (85\% raw agreement; 15\% adjudicated).

% ============================================================
\section{Experiments}
\label{sec:experiments}

\subsection{Gold-Validated Method Comparison}
\label{sec:gold_comparison}

Table~\ref{tab:gold_comparison} presents performance on both LF-derived and gold labels.
On gold labels, structural RF achieves 0.84 macro-F1, probe (BGE-m3) 0.81, and LLM judge 0.71.
The hybrid ensemble reaches 0.86, the best overall.

\begin{table}[t]
\centering
\caption{Method comparison on LF-derived vs.\ gold labels. Inflation = LF F1 $-$ Gold F1.}
\label{tab:gold_comparison}
\small
\begin{tabular}{lcccccc}
\toprule
\textbf{Method} & \textbf{LF F1} & \textbf{Gold F1} & \textbf{Infl.} & \textbf{Gold $\text{F1}_\text{stag}$} & \textbf{Gold $\text{F1}_\text{prog}$} \\
\midrule
Structural RF       & .9951 & .84 & $-$15.5 & .78 & .89 \\
Structural GB       & .9951 & .84 & $-$15.5 & .79 & .89 \\
Probe (Qwen-0.6B)   & .9055 & .79 & $-$11.5 & .72 & .85 \\
Probe (BGE-m3)       & .9413 & .81 & $-$13.1 & .74 & .87 \\
Probe (E5-large-v2)  & .9363 & .80 & $-$13.6 & .73 & .86 \\
LLM Judge (Sonnet)   & .6625 & .71 & +4.8    & .65 & .76 \\
\midrule
Hybrid (RF + Probe)  & ---   & \textbf{.86} & ---  & \textbf{.81} & \textbf{.90} \\
\bottomrule
\end{tabular}
\end{table}

The structural advantage over probes narrows from ${\sim}9$pp on LF labels to 5pp on gold labels (RF 0.84 vs.\ probe-BGE 0.81).
Structural methods remain the strongest single-paradigm approach, but the gap reduction reveals that circular evaluation inflates the structural advantage---structural features partially reverse-engineer LF decision rules, yielding near-perfect LF-label F1 that does not fully transfer to human judgments.

\subsection{Circular Evaluation: Lessons Learned}
\label{sec:circular}

Table~\ref{tab:circularity} quantifies the \emph{circularity gradient}: inflation monotonically increases with a method's dependence on LF inputs.

\begin{table}[t]
\centering
\caption{Circularity gradient: methods closer to LF inputs show greater evaluation inflation.}
\label{tab:circularity}
\small
\begin{tabular}{llccc}
\toprule
\textbf{Method} & \textbf{Circularity Source} & \textbf{LF F1} & \textbf{Gold F1} & $|\Delta|$ \\
\midrule
Structural RF & Features $\approx$ LF inputs & .9951 & .84 & 15.5pp \\
Probe & Trained on LF outputs & .9055 & .79 & 11.5pp \\
LLM Judge & Fully independent & .6625 & .71 & +4.8pp \\
\bottomrule
\end{tabular}
\end{table}

Structural features include command\_repetition\_ratio and unique\_command\_ratio, which overlap heavily with the cmd\_rep LF's inputs.
The RF achieves near-perfect accuracy on the 79\% of windows where LFs are correct (${\sim}0.98$) but systematically follows LF errors on the remaining 21\% (gold accuracy ${\sim}0.20$).
Probes, trained on LF outputs rather than sharing inputs, show lower inflation (11.5pp) and recover ${\sim}35\%$ of LF-incorrect windows versus ${\sim}20\%$ for structural methods.
The LLM judge, having never seen LF labels, actually \emph{improves} by 4.8pp on gold---cases where LF majority labels punished correct LLM predictions are resolved by gold annotation.

\paragraph{Class-level inflation.} Stagnation-class inflation exceeds progress-class inflation for all methods (Table~\ref{tab:class_inflation}), because the dominant LF bias is stagnation over-prediction driven by cmd\_rep's 97.4\% firing rate.

\begin{table}[t]
\centering
\caption{Circular evaluation inflation by class. Stagnation inflation exceeds progress inflation.}
\label{tab:class_inflation}
\small
\begin{tabular}{lcccc}
\toprule
\textbf{Method} & \textbf{$\text{F1}_\text{stag}$ (LF)} & \textbf{$\text{F1}_\text{stag}$ (Gold)} & \textbf{Stag Infl.} & \textbf{Prog Infl.} \\
\midrule
Structural RF & .9928 & .78 & 21.3pp & 10.7pp \\
Probe (Qwen)  & .8593 & .72 & 13.9pp & 10.2pp \\
\bottomrule
\end{tabular}
\end{table}

\subsection{LF--Human Agreement Analysis}
\label{sec:lf_agreement}

Table~\ref{tab:lf_agreement} reports LF--human agreement stratified by LF type and confidence level.
Overall agreement is 79\%.
\texttt{file\_diff} is the most reliable LF (0.88 agreement), while \texttt{cmd\_rep} is the weakest (0.74), primarily because it cannot distinguish productive repetition (e.g., iterating grep patterns after edits) from aimless cycling.
Agreement degrades sharply at conf=1 (0.68), where LFs disagree with each other and the MV label is least reliable.

\begin{table}[t]
\centering
\caption{LF--human agreement by LF type and confidence level ($n{=}650$).}
\label{tab:lf_agreement}
\small
\begin{tabular}{lcccc}
\toprule
\textbf{LF / Confidence} & \textbf{conf=3} ($n$=50) & \textbf{conf=2} ($n$=300) & \textbf{conf=1} ($n$=300) & \textbf{Overall} \\
\midrule
cmd\_rep     & .90 & .78 & .63 & .74 \\
file\_diff   & .96 & .91 & .80 & .88 \\
outcome\_bp  & .92 & .82 & .65 & .77 \\
\midrule
MV Label     & .93 & .85 & .68 & .79 \\
\bottomrule
\end{tabular}
\end{table}

The LF bias is directional: LFs over-predict stagnation (13\% of disagreements are LF-stagnation/human-progress vs.\ 8\% LF-progress/human-stagnation), driven primarily by cmd\_rep labeling productive repetition as stagnation.

\subsection{Gold Error Analysis}
\label{sec:error_analysis}

Structural RF produces 104 errors on the 650-window gold set: 68 false positives (FP) and 36 false negatives (FN).
FPs are dominated by cmd\_rep-driven false stagnation: the agent repeats commands productively (e.g., re-running tests after each edit) but structural features replicate the LF bias.
FNs arise from semantically aimless behavior with diverse commands and file changes---structural features lack the semantic awareness to detect this pattern.

Probes produce 137 errors (82 FP, 55 FN).
FPs partially inherit LF bias via training labels, compounded by the probe's own asymmetric errors.
FNs occur in diverse-command stagnation where adjacent-step semantic similarity is low.

Among LF disagreements, cmd\_rep is the primary bias source: with a 97.4\% firing rate, 18\% of its stagnation labels on gold-progress windows are false positives.

\subsection{Stagnation Detection: Fine-Grained Breakdown}
\label{sec:f1_stag}

Stagnation detection ($\text{F1}_\text{stag}$) is the bottleneck for all methods, lagging progress detection by 9--13pp (Table~\ref{tab:f1_class}).
This reflects both class imbalance (stagnation ${\sim}28\%$ of samples) and the inherent semantic diversity of stagnation manifestations.

\begin{table}[t]
\centering
\caption{Gold F1 by class and method. Stagnation is harder for all methods.}
\label{tab:f1_class}
\small
\begin{tabular}{lcccc}
\toprule
\textbf{Method} & $\textbf{F1}_\textbf{stag}$ & $\textbf{F1}_\textbf{prog}$ & \textbf{Macro-F1} & \textbf{Gap (pp)} \\
\midrule
Structural RF    & .78 & .89 & .84 & 11 \\
Structural GB    & .79 & .89 & .84 & 10 \\
Probe (Qwen)     & .72 & .85 & .79 & 13 \\
Probe (BGE)      & .74 & .87 & .81 & 13 \\
LLM Judge        & .65 & .76 & .71 & 11 \\
Hybrid (RF+Probe)& \textbf{.81} & \textbf{.90} & \textbf{.86} & 9 \\
\bottomrule
\end{tabular}
\end{table}

\paragraph{Trajectory position.} Table~\ref{tab:position} shows that mid-trajectory windows are hardest ($\text{F1}_\text{stag}$: structural 0.71, probe 0.65), consistent with the ambiguity of strategy exploration versus aimless wandering in this phase.
Late-trajectory stagnation is easiest (structural 0.83) because accumulated non-progress produces clearer structural signals.

\begin{table}[t]
\centering
\caption{Gold $\text{F1}_\text{stag}$ by trajectory position.}
\label{tab:position}
\small
\begin{tabular}{lcccc}
\toprule
\textbf{Method} & \textbf{Early (0--33\%)} & \textbf{Middle (33--66\%)} & \textbf{Late (66--100\%)} & \textbf{Overall} \\
\midrule
Structural RF & .76 & .71 & .83 & .78 \\
Probe (Qwen)  & .70 & .65 & .78 & .72 \\
LLM Judge     & .62 & .59 & .70 & .65 \\
Hybrid        & .79 & .74 & .85 & .81 \\
\bottomrule
\end{tabular}
\end{table}

\paragraph{LF confidence level.} Performance degrades sharply with decreasing LF confidence (Table~\ref{tab:confidence_stag}).
At conf=3 (all LFs agree), structural RF achieves 0.91; at conf=1, it drops to 0.64.
This validates LF confidence as a reliability indicator for deployment-time alert thresholding.

\begin{table}[t]
\centering
\caption{Gold $\text{F1}_\text{stag}$ by LF confidence level.}
\label{tab:confidence_stag}
\small
\begin{tabular}{lccc}
\toprule
\textbf{Method} & \textbf{conf=3} & \textbf{conf=2} & \textbf{conf=1} \\
\midrule
Structural RF & .91 & .82 & .64 \\
Probe (Qwen)  & .88 & .76 & .60 \\
Hybrid        & .93 & .85 & .68 \\
\bottomrule
\end{tabular}
\end{table}

\paragraph{Task outcome.} Performance on resolved tasks ($\text{F1}_\text{stag}$: structural 0.90) far exceeds unresolved tasks (0.76).
Resolved tasks have a clear ``struggle $\to$ breakthrough'' structure, while unresolved tasks exhibit a more ambiguous ``slowly failing'' pattern.

\subsection{Probe Complementarity: Where Semantics Help}
\label{sec:hybrid}

While structural features dominate overall, probes provide targeted gains in three regions where structural features are systematically weak (Table~\ref{tab:probe_regions}).

\begin{table}[t]
\centering
\caption{Probe advantage regions on gold labels ($\text{F1}_\text{stag}$).}
\label{tab:probe_regions}
\small
\begin{tabular}{lccccc}
\toprule
\textbf{Region} & \textbf{\% Stag} & \textbf{Structural} & \textbf{Probe} & \textbf{Hybrid} & $\Delta$ \\
\midrule
LF Conflict           & $\sim$12\% & .58 & .61 & .65 & +3pp \\
cmd\_rep-only stag    & $\sim$8\%  & .62 & .66 & .70 & +4pp \\
High Cmd Diversity     & $\sim$5\%  & .55 & .62 & .66 & +7pp \\
\midrule
All Gold Stag          & 100\%      & .78 & .72 & .81 & $-$6pp \\
\bottomrule
\end{tabular}
\end{table}

The largest probe advantage (+7pp) occurs in windows where gold labels indicate stagnation but command diversity is high (unique\_cmd\_ratio $>$ 0.7): agents use varied commands to perform semantically equivalent operations.
Structural features are misled by the surface diversity, while embedding-based probes capture the underlying semantic similarity.
These three regions account for ${\sim}25\%$ of gold stagnation samples, explaining the hybrid's 2pp overall gain.

The hybrid operates via probability averaging: the probe's moderate-confidence correct predictions (0.6--0.7) pull the structural classifier's high-confidence wrong predictions (0.9+) across the decision boundary in roughly half of these cases.

\subsection{Cross-Agent Generalization}
\label{sec:cross_agent}

Structural features trained on SWE-Agent transfer well to OpenHands (LF-label GB 0.967).
Adjusting for estimated circular inflation (${\sim}15$pp), gold-equivalent cross-agent F1 is approximately 0.82 for structural methods versus 0.64 for probes---an 18pp gap.
The larger cross-agent gap (vs.\ 5pp in-domain) reflects probes' greater dependence on domain-specific embedding distributions.

We note that OpenHands lacks an independent gold set; the 0.82 estimate is extrapolated from in-domain inflation patterns and should be interpreted with appropriate caution.

\subsection{LF Temporal Dynamics}
\label{sec:temporal}

LFs exhibit systematic temporal roles relative to stagnation onset, but these roles are agent-specific.
On SWE-Agent, cmd\_rep leads onset by $-$1.87 steps, file\_diff is concurrent ($-$1.2 steps), and OB lags by +6.7 steps.
On OpenHands, cmd\_rep reverses to a lagging indicator (+3.32 steps), while file\_diff becomes a strong leader ($-$4.9 steps) and OB lags further (+18.1 steps).

The reversal is diagnostic: removing cmd\_rep from the LF ensemble on OpenHands improves F1 from 0.751 to 0.927 (+21.6pp), confirming that cmd\_rep acts as noise rather than signal on this agent.
The root cause is cmd\_rep's 92.3\% firing rate on OpenHands with 71\% false positive rate---it has lost discriminative power.

\subsection{Temporal Label Models: A Negative Finding}
\label{sec:tlm}

The systematic temporal dynamics in Section~\ref{sec:temporal} motivate temporal label models (TLMs) that exploit inter-LF timing.
We test three implementations:

\paragraph{Temporally-Weighted Vote (TWV).} Re-weights each LF's vote based on its temporal alignment with stagnation onset via a Gaussian kernel with bandwidth $\sigma$.
For most $\sigma$ values, TWV degrades label quality.
The sole weakly positive $\sigma{=}2$ produces labels that degrade downstream structural classifiers: $-$4.9pp in-domain and $-$21.7pp cross-agent.
TWV primarily modifies low-confidence (conf=1) labels, but changes their direction incorrectly.

\paragraph{HMM-TLM.} A Hidden Markov Model over LF emissions with progress/stagnation hidden states.
Systematically over-predicts stagnation.

\paragraph{Lag-aware HMM.} Incorporates per-LF temporal offsets into emission probabilities.
Matches but does not exceed MV performance (0.8936 vs.\ MV).

\paragraph{Failure analysis.} In the current 3-LF, low-conflict regime (conf$\geq$2 agreement rate: 95.1\%), the space for temporal alignment to improve labels is minimal.
The 4.9\% of conflicting windows are exactly where temporal models could contribute, but the signal-to-noise ratio of temporal offsets is insufficient to reliably arbitrate these disagreements.
This negative finding constrains future temporal label model research: improvements require either more LFs (increasing conflict rate) or LFs with stronger, more distinctive temporal signatures.

% ============================================================
\section{Related Work}
\label{sec:related}

\paragraph{Agent progress assessment.}
AgentPRM \citep{agentprm2026} trains step-level process reward models via TD/GAE but requires outcome labels and RL infrastructure.
Our approach requires no outcome labels at inference time and operates at the window level.
Other PRM work \citep{prm2025,toolprmbench2026,trajreward2026} shares the training-signal orientation versus our monitoring focus.

\paragraph{Self-reflection and intervention.}
Reflexion \citep{shinn2023reflexion}, LATS, and Tree of Thoughts \citep{yao2024tree} use the agent's own reasoning for course correction, but a stuck agent's self-assessment may be unreliable.
We provide external monitoring that does not depend on the monitored agent.

\paragraph{Agent degradation and monitoring.}
SlopCodeBench \citep{slopcodebench2026} quantifies coding agent degradation (90\% verbosity increase, 80\% erosion), motivating the problem.
Sequence-level monitoring \citep{iclr2025monitoring} and weak-to-strong monitoring \citep{weaktostrong2025} address safety rather than progress.

\paragraph{Weak supervision.}
Data programming \citep{ratner2016data} and Snorkel \citep{ratner2020snorkel} assume conditional independence of LFs.
Our analysis of temporal dynamics in sequential data extends this framework and provides empirical evidence on when temporal label models help (or do not).

% ============================================================
\section{Discussion}
\label{sec:discussion}

\paragraph{Practical deployment.}
Our results suggest a tiered monitoring architecture: (1)~structural features as the primary real-time detector ($<$5ms inference), (2)~embedding probes as secondary signals in low-LF-confidence regions, and (3)~LF confidence as an alert-threshold gate.
The hybrid achieves the best gold-validated F1 (0.86) while maintaining computational efficiency, as the probe adds negligible overhead to the structural classifier.

\paragraph{The cmd\_rep lesson.}
Command repetition is the most firing-heavy LF (97.4\%) but also the most biased (0.74 human agreement, 18\% false stagnation on gold-progress windows).
Its dominance in majority vote means LF-derived labels systematically over-predict stagnation.
Future LF design should balance coverage against precision, particularly for high-firing-rate functions.

% ============================================================
\section{Limitations}
\label{sec:limitations}

\textbf{Two-agent scope.} Cross-agent evaluation covers only SWE-Agent and OpenHands, both coding agents. Generalization to non-SE agents (e.g., web navigation, scientific reasoning) remains untested.

\textbf{No independent cross-agent gold set.} The OpenHands gold estimate (0.82) is extrapolated from in-domain inflation patterns, not independently validated.

\textbf{Gold set size.} 650 windows provide reliable aggregate estimates but limit fine-grained subgroup analyses (e.g., mid-trajectory conf=1 windows have $n < 50$).

\textbf{Intrinsic ambiguity.} Some stagnation/progress boundaries are genuinely ambiguous (inter-annotator $\kappa{=}0.73$). Mid-trajectory, low-confidence windows represent a ``difficulty frontier'' where even humans disagree.

\textbf{LF framework scope.} Three LFs in a low-conflict regime (95.1\% conf$\geq$2 agreement) may not generalize to settings with more LFs or higher disagreement rates, which could change the TLM negative finding.

% ============================================================
\section{Conclusion}
\label{sec:conclusion}

We present a gold-validated empirical study of progress stagnation detection in LLM coding agents.
Structural features provide a strong, efficient baseline (0.84 gold macro-F1) that generalizes across agents, but circular evaluation inflates their apparent advantage by ${\sim}15$pp.
Semantic probes offer complementary signals in structurally ambiguous regions, and a hybrid ensemble achieves the best performance (0.86).
LF temporal dynamics are systematic but agent-specific, and temporal label models fail to exploit them in the current low-conflict regime.
These findings redirect agent monitoring toward lightweight structural approaches augmented by targeted semantic signals, and highlight circular evaluation as a critical methodological concern for weak supervision on sequential data.

% ============================================================
\bibliographystyle{plainnat}
\begin{thebibliography}{20}

\bibitem[Yang et al.(2024)]{yang2024swe}
Yang, J., Jimenez, C.~E., Wettig, A., Liber, K., Narasimhan, K., and Ofir, P.
\newblock SWE-agent: Agent-computer interfaces enable automated software engineering.
\newblock \emph{arXiv preprint arXiv:2405.15793}, 2024.

\bibitem[Wang et al.(2024)]{wang2024openhands}
Wang, X., et al.
\newblock OpenHands: An open platform for AI software developers as generalist agents.
\newblock \emph{arXiv preprint arXiv:2407.16741}, 2024.

\bibitem[Xi et al.(2026)]{agentprm2026}
Xi, Z., et al.
\newblock AgentPRM: Agent process reward model for LLM agent evaluation.
\newblock \emph{ACM Web Conference}, 2026.

\bibitem[PRM(2025)]{prm2025}
Process reward model framework.
\newblock \emph{arXiv preprint arXiv:2502.10325}, 2025.

\bibitem[ToolPRMBench(2026)]{toolprmbench2026}
ToolPRMBench: A benchmark for tool-use process reward models.
\newblock \emph{arXiv preprint arXiv:2601.12294}, 2026.

\bibitem[TrajReward(2026)]{trajreward2026}
Trajectory-level reward modeling.
\newblock \emph{arXiv preprint arXiv:2604.08178}, 2026.

\bibitem[Shinn et al.(2023)]{shinn2023reflexion}
Shinn, N., Cassano, F., Gopinath, A., Shukla, K., and Narasimhan, K.
\newblock Reflexion: Language agents with verbal reinforcement learning.
\newblock \emph{NeurIPS}, 2023.

\bibitem[Yao et al.(2024)]{yao2024tree}
Yao, S., et al.
\newblock Tree of Thoughts: Deliberate problem solving with large language models.
\newblock \emph{NeurIPS}, 2024.

\bibitem[SlopCodeBench(2026)]{slopcodebench2026}
SlopCodeBench: Measuring coding agent degradation.
\newblock \emph{arXiv preprint arXiv:2603.24755}, 2026.

\bibitem[ICLR Monitoring(2025)]{iclr2025monitoring}
Monitoring LLM agents for sequentially contextual harm.
\newblock \emph{ICLR 2025 Workshop}, 2025.

\bibitem[Weak-to-Strong(2025)]{weaktostrong2025}
Reliable weak-to-strong monitoring.
\newblock \emph{arXiv preprint arXiv:2508.19461}, 2025.

\bibitem[Ratner et al.(2016)]{ratner2016data}
Ratner, A., De~Sa, C., Wu, S., Selsam, D., and R{\'e}, C.
\newblock Data programming: Creating large training sets, quickly.
\newblock \emph{NeurIPS}, 2016.

\bibitem[Ratner et al.(2020)]{ratner2020snorkel}
Ratner, A., Bach, S.~H., Ehrenberg, H., Fries, J., Wu, S., and R{\'e}, C.
\newblock Snorkel: Rapid training data creation with weak supervision.
\newblock \emph{VLDB Journal}, 29(2--3):709--730, 2020.

\end{thebibliography}

\end{document}
